\begin{abstract}
Keystroke signals are attractive triggers for timely help if they change before a learner stops producing text. We
tested whether the 30-second change, at the start of a pause, in two signals of a keystroke state estimator with
hand-set coefficients (the probability of its Stuck state and a smoothed friction axis) distinguishes pauses that became
stalls in character production from matched pauses of 2--30~s that did not. The test was preregistered on a random
holdout of an introductory programming dataset, with an equivalence band of $[0.45, 0.55]$ for the AUC set in advance.
Because a leak had revealed the direction on seen writers, it is a confirmation, not a test of a prediction that could
fail. In the holdout (139 students), the AUC fell within the band for both signals, slightly below 0.5 (0.489 and 0.487;
the false-equivalence rate is at most 5.2\% [Wilson 95\% interval 3.6--7.5\%] among the configurations examined in the
calibration on seen writers, and is not an upper bound over the whole null). The Stuck-state probability is mostly at
the floor of the estimate, and claims about run-up change rely on the friction axis and on the pairs not at the floor.
In exploration on seen writers, only injected changes far larger than those observed between stalls and pauses moved
the friction axis outside the band. On the same pairs, learners on standard keystroke quantities exceeded both signals
by 0.05--0.08 without excluding the upper side of the band. On this input, the 30-second change gave no basis for a
trigger; the evaluation code is released for checking estimators of this type against a band set in advance.
Preregistration: \url{https://osf.io/mrkxt}; results: \url{https://osf.io/uj6cy}; code: [archive to be added before
submission].\\

{\parindent0pt
\textbf{Keywords:} keystroke logging, stalls, equivalence testing, preregistration, matched controls, introductory
programming
}
\end{abstract}
